%% ECON 282E -- Session 5, Deck A2: The Model, and the Parameter That Scales Risk
%% Source: Quant_Macro Lec.9 frames 9-15, plus the two "Taylor's theorem" frames
%% that the successor deck dropped (Lec5-7_Dynamic_models.tex:1608, 1624).
%% Exposition follows Fernandez-Villaverde & Guerron-Quintana, Solution Methods V.
%% Numbers from tools/figures/s05_numbers.json, key "brock_mirman".
%% NOTATION: the perturbation parameter is chi.  The literature writes lambda;
%% this course reserves lambda for the invariant distribution (CLAUDE.md).

%% ECON 282E -- Foundations of Macroeconomics (UC Riverside, Fall 2026)
%% Shared Beamer preamble for every deck in slides/S01 ... slides/S10.
%%
%% Usage (a deck lives in slides/S0X/ and is compiled from that folder):
%%   %% ECON 282E -- Foundations of Macroeconomics (UC Riverside, Fall 2026)
%% Shared Beamer preamble for every deck in slides/S01 ... slides/S10.
%%
%% Usage (a deck lives in slides/S0X/ and is compiled from that folder):
%%   %% ECON 282E -- Foundations of Macroeconomics (UC Riverside, Fall 2026)
%% Shared Beamer preamble for every deck in slides/S01 ... slides/S10.
%%
%% Usage (a deck lives in slides/S0X/ and is compiled from that folder):
%%   \input{../preamble/course_preamble.tex}
%%   \decktitle{1}{A1}{Who I Am \& Why This Course}{September 24, 2026}
%%   \begin{document}
%%   \makedecktitle
%%   ...
%%
%% Adapted from Courses/AI-ECON-2026/source/shared_preamble.tex (Metropolis theme,
%% concept/caution/example/takeaway boxes, \sectiondivider). Changes for this course:
%%   * course title block (\decktitle / \makedecktitle) instead of \makebeamertitle
%%   * \zh{...} typesets Chinese (book titles) under pdflatex via CJKutf8
%%   * researchbox for the "Research in the wild" frames
%%   * section pages off; TikZ libraries and the course map loaded here
%% Build with pdflatex only (two passes); tools/build_slides.py does this for you.

\documentclass[10pt,english,aspectratio=169]{beamer}

\usepackage[T1]{fontenc}
\usepackage{textcomp}
\usepackage[utf8]{inputenc}
\usepackage{babel}
\usepackage{CJKutf8}

\usepackage{amsmath, amssymb, amsthm, graphicx}
\usepackage{booktabs, multirow, tabularx}
\usepackage{tikz}
\usetikzlibrary{positioning,arrows.meta,shapes,calc,fit,backgrounds}
\usepackage{pgfplots}
\pgfplotsset{compat=1.16}
\usepackage[most]{tcolorbox}
\tcbuselibrary{skins, breakable}
\usepackage{listings}
\usepackage{xcolor}

\lstset{
  basicstyle=\ttfamily\small,
  keywordstyle=\color{blue!80!black}\bfseries,
  commentstyle=\color{green!50!black}\itshape,
  stringstyle=\color{orange!80!black},
  backgroundcolor=\color{gray!8},
  frame=single,
  rulecolor=\color{gray!40},
  breaklines=true,
  showstringspaces=false,
  numbers=none,
}

\usetheme{metropolis}
\usefonttheme{professionalfonts}
\metroset{sectionpage=none}

\definecolor{LightGray}{RGB}{230,230,230}
\definecolor{RedTitle}{RGB}{0,0,200}
\definecolor{VeryLightGray}{RGB}{245,245,245}
\definecolor{DarkText}{RGB}{0,0,0}
\definecolor{UCRBlue}{RGB}{0,45,106}
\definecolor{UCRGold}{RGB}{241,171,0}

% Stage colours of the course arc (used by the course map and elsewhere)
\colorlet{StageTrunk}{blue!12}
\colorlet{StageBranch}{cyan!14}
\colorlet{StageMachine}{gray!18}
\colorlet{StageWall}{red!14}
\colorlet{StageTools}{orange!20}
\colorlet{StageData}{green!16}
\colorlet{StageAgents}{violet!14}

\hypersetup{bookmarks=false,breaklinks=true,pdfborder={0 0 0},colorlinks=true,
  linkcolor=DarkText,urlcolor=RedTitle}

\setbeamercolor{background canvas}{bg=white}
\setbeamercolor{title}{fg=RedTitle}
\setbeamercolor{frametitle}{bg=VeryLightGray,fg=DarkText}
\setbeamercolor{normal text}{fg=DarkText}
\setbeamercolor{alerted text}{fg=RedTitle}
\setbeamersize{text margin left=5mm, text margin right=5mm}
\addtobeamertemplate{frametitle}{}{\vspace{-0.5em}}

\setbeamertemplate{navigation symbols}{}
\setbeamertemplate{footline}{%
  \hspace{5mm}{\usebeamerfont{page number in head/foot}\color{black!45}%
  ECON 282E \,$\cdot$\, \insertshortdeck}\hfill%
  \usebeamercolor[fg]{page number in head/foot}%
  \usebeamerfont{page number in head/foot}%
  \insertframenumber\,/\,\inserttotalframenumber\kern1em\vskip2pt%
}

% ---------------------------------------------------------------------------
% Chinese text under pdflatex (book titles etc.)
\newcommand{\zh}[1]{\begin{CJK*}{UTF8}{gbsn}#1\end{CJK*}}

% ---------------------------------------------------------------------------
% Deck title block
%   \decktitle{session}{deck id}{title}{date}
\newcommand{\insertshortdeck}{}
\newcommand{\decksession}{}
\newcommand{\deckid}{}
\newcommand{\decktitletext}{}
\newcommand{\deckdate}{}
\newcommand{\decktitle}[4]{%
  \renewcommand{\decksession}{#1}%
  \renewcommand{\deckid}{#2}%
  \renewcommand{\decktitletext}{#3}%
  \renewcommand{\deckdate}{#4}%
  \renewcommand{\insertshortdeck}{Session #1 \,$\cdot$\, Deck #1.#2}%
  \title{#3}%
  \author{Zhigang Feng}%
  \date{#4}%
}
\newcommand{\makedecktitle}{%
  \begin{frame}[plain,noframenumbering]
    \begin{tikzpicture}[remember picture,overlay]
      \fill[UCRBlue] (current page.north west) rectangle ([yshift=-0.55cm]current page.north east);
      \fill[UCRGold] ([yshift=-0.55cm]current page.north west) rectangle ([yshift=-0.62cm]current page.north east);
      \node[anchor=west, text=white, font=\small\bfseries] at ([xshift=5mm,yshift=-0.275cm]current page.north west)
        {ECON 282E \,$\cdot$\, Foundations of Macroeconomics};
      \node[anchor=east, text=white, font=\small] at ([xshift=-5mm,yshift=-0.275cm]current page.north east)
        {UC Riverside \,$\cdot$\, Fall 2026};
    \end{tikzpicture}
    \vspace*{1.1cm}
    {\usebeamerfont{title}\color{black!55}\large Session \decksession\ \,$\cdot$\, Deck \decksession.\deckid\par}
    \vspace{0.25cm}
    {\usebeamerfont{title}\color{RedTitle}\LARGE\bfseries \decktitletext\par}
    \vspace{0.7cm}
    {\large Zhigang Feng\par}
    \vspace{0.15cm}
    {\color{black!65}\deckdate\par}
    \vfill
    {\footnotesize\itshape\color{black!55}Build it on paper $\;\to\;$ solve it on the machine $\;\to\;$ push it to the frontier with AI\par}
  \end{frame}}

% ---------------------------------------------------------------------------
% Section divider (plain frame, not counted)
\newcommand{\sectiondivider}[2]{%
\begin{frame}[plain,noframenumbering]
\begin{center}
\vfill
{\Large\textbf{#1}\par}
\vspace{0.35cm}
{\normalsize #2\par}
\vfill
\end{center}
\end{frame}
}

% ---------------------------------------------------------------------------
% Boxes
\newtcolorbox{conceptbox}[1][]{colback=blue!3, colframe=RedTitle!55, boxrule=0.35mm,
  arc=1mm, left=2mm, right=2mm, top=1mm, bottom=1mm, #1}
\newtcolorbox{cautionbox}[1][]{colback=red!3, colframe=red!50!black, boxrule=0.35mm,
  arc=1mm, left=2mm, right=2mm, top=1mm, bottom=1mm, #1}
\newtcolorbox{examplebox}[1][]{colback=green!3, colframe=green!45!black, boxrule=0.35mm,
  arc=1mm, left=2mm, right=2mm, top=1mm, bottom=1mm, #1}
\newtcolorbox{takeawaybox}[1][]{colback=VeryLightGray, colframe=RedTitle!55, boxrule=0.35mm,
  arc=1mm, left=2mm, right=2mm, top=1mm, bottom=1mm, #1}
\newtcolorbox{researchbox}[1][]{enhanced, colback=violet!4, colframe=violet!60!black,
  boxrule=0.35mm, arc=1mm, left=2mm, right=2mm, top=1mm, bottom=1mm,
  fonttitle=\bfseries\small, title={Research in the wild}, #1}

\newcommand{\compacttables}{\renewcommand{\arraystretch}{1.15}}
\newcommand{\src}[1]{{\tiny\color{black!55}#1}}

% ---------------------------------------------------------------------------
\input{../preamble/course_map.tex}

%%   \decktitle{1}{A1}{Who I Am \& Why This Course}{September 24, 2026}
%%   \begin{document}
%%   \makedecktitle
%%   ...
%%
%% Adapted from Courses/AI-ECON-2026/source/shared_preamble.tex (Metropolis theme,
%% concept/caution/example/takeaway boxes, \sectiondivider). Changes for this course:
%%   * course title block (\decktitle / \makedecktitle) instead of \makebeamertitle
%%   * \zh{...} typesets Chinese (book titles) under pdflatex via CJKutf8
%%   * researchbox for the "Research in the wild" frames
%%   * section pages off; TikZ libraries and the course map loaded here
%% Build with pdflatex only (two passes); tools/build_slides.py does this for you.

\documentclass[10pt,english,aspectratio=169]{beamer}

\usepackage[T1]{fontenc}
\usepackage{textcomp}
\usepackage[utf8]{inputenc}
\usepackage{babel}
\usepackage{CJKutf8}

\usepackage{amsmath, amssymb, amsthm, graphicx}
\usepackage{booktabs, multirow, tabularx}
\usepackage{tikz}
\usetikzlibrary{positioning,arrows.meta,shapes,calc,fit,backgrounds}
\usepackage{pgfplots}
\pgfplotsset{compat=1.16}
\usepackage[most]{tcolorbox}
\tcbuselibrary{skins, breakable}
\usepackage{listings}
\usepackage{xcolor}

\lstset{
  basicstyle=\ttfamily\small,
  keywordstyle=\color{blue!80!black}\bfseries,
  commentstyle=\color{green!50!black}\itshape,
  stringstyle=\color{orange!80!black},
  backgroundcolor=\color{gray!8},
  frame=single,
  rulecolor=\color{gray!40},
  breaklines=true,
  showstringspaces=false,
  numbers=none,
}

\usetheme{metropolis}
\usefonttheme{professionalfonts}
\metroset{sectionpage=none}

\definecolor{LightGray}{RGB}{230,230,230}
\definecolor{RedTitle}{RGB}{0,0,200}
\definecolor{VeryLightGray}{RGB}{245,245,245}
\definecolor{DarkText}{RGB}{0,0,0}
\definecolor{UCRBlue}{RGB}{0,45,106}
\definecolor{UCRGold}{RGB}{241,171,0}

% Stage colours of the course arc (used by the course map and elsewhere)
\colorlet{StageTrunk}{blue!12}
\colorlet{StageBranch}{cyan!14}
\colorlet{StageMachine}{gray!18}
\colorlet{StageWall}{red!14}
\colorlet{StageTools}{orange!20}
\colorlet{StageData}{green!16}
\colorlet{StageAgents}{violet!14}

\hypersetup{bookmarks=false,breaklinks=true,pdfborder={0 0 0},colorlinks=true,
  linkcolor=DarkText,urlcolor=RedTitle}

\setbeamercolor{background canvas}{bg=white}
\setbeamercolor{title}{fg=RedTitle}
\setbeamercolor{frametitle}{bg=VeryLightGray,fg=DarkText}
\setbeamercolor{normal text}{fg=DarkText}
\setbeamercolor{alerted text}{fg=RedTitle}
\setbeamersize{text margin left=5mm, text margin right=5mm}
\addtobeamertemplate{frametitle}{}{\vspace{-0.5em}}

\setbeamertemplate{navigation symbols}{}
\setbeamertemplate{footline}{%
  \hspace{5mm}{\usebeamerfont{page number in head/foot}\color{black!45}%
  ECON 282E \,$\cdot$\, \insertshortdeck}\hfill%
  \usebeamercolor[fg]{page number in head/foot}%
  \usebeamerfont{page number in head/foot}%
  \insertframenumber\,/\,\inserttotalframenumber\kern1em\vskip2pt%
}

% ---------------------------------------------------------------------------
% Chinese text under pdflatex (book titles etc.)
\newcommand{\zh}[1]{\begin{CJK*}{UTF8}{gbsn}#1\end{CJK*}}

% ---------------------------------------------------------------------------
% Deck title block
%   \decktitle{session}{deck id}{title}{date}
\newcommand{\insertshortdeck}{}
\newcommand{\decksession}{}
\newcommand{\deckid}{}
\newcommand{\decktitletext}{}
\newcommand{\deckdate}{}
\newcommand{\decktitle}[4]{%
  \renewcommand{\decksession}{#1}%
  \renewcommand{\deckid}{#2}%
  \renewcommand{\decktitletext}{#3}%
  \renewcommand{\deckdate}{#4}%
  \renewcommand{\insertshortdeck}{Session #1 \,$\cdot$\, Deck #1.#2}%
  \title{#3}%
  \author{Zhigang Feng}%
  \date{#4}%
}
\newcommand{\makedecktitle}{%
  \begin{frame}[plain,noframenumbering]
    \begin{tikzpicture}[remember picture,overlay]
      \fill[UCRBlue] (current page.north west) rectangle ([yshift=-0.55cm]current page.north east);
      \fill[UCRGold] ([yshift=-0.55cm]current page.north west) rectangle ([yshift=-0.62cm]current page.north east);
      \node[anchor=west, text=white, font=\small\bfseries] at ([xshift=5mm,yshift=-0.275cm]current page.north west)
        {ECON 282E \,$\cdot$\, Foundations of Macroeconomics};
      \node[anchor=east, text=white, font=\small] at ([xshift=-5mm,yshift=-0.275cm]current page.north east)
        {UC Riverside \,$\cdot$\, Fall 2026};
    \end{tikzpicture}
    \vspace*{1.1cm}
    {\usebeamerfont{title}\color{black!55}\large Session \decksession\ \,$\cdot$\, Deck \decksession.\deckid\par}
    \vspace{0.25cm}
    {\usebeamerfont{title}\color{RedTitle}\LARGE\bfseries \decktitletext\par}
    \vspace{0.7cm}
    {\large Zhigang Feng\par}
    \vspace{0.15cm}
    {\color{black!65}\deckdate\par}
    \vfill
    {\footnotesize\itshape\color{black!55}Build it on paper $\;\to\;$ solve it on the machine $\;\to\;$ push it to the frontier with AI\par}
  \end{frame}}

% ---------------------------------------------------------------------------
% Section divider (plain frame, not counted)
\newcommand{\sectiondivider}[2]{%
\begin{frame}[plain,noframenumbering]
\begin{center}
\vfill
{\Large\textbf{#1}\par}
\vspace{0.35cm}
{\normalsize #2\par}
\vfill
\end{center}
\end{frame}
}

% ---------------------------------------------------------------------------
% Boxes
\newtcolorbox{conceptbox}[1][]{colback=blue!3, colframe=RedTitle!55, boxrule=0.35mm,
  arc=1mm, left=2mm, right=2mm, top=1mm, bottom=1mm, #1}
\newtcolorbox{cautionbox}[1][]{colback=red!3, colframe=red!50!black, boxrule=0.35mm,
  arc=1mm, left=2mm, right=2mm, top=1mm, bottom=1mm, #1}
\newtcolorbox{examplebox}[1][]{colback=green!3, colframe=green!45!black, boxrule=0.35mm,
  arc=1mm, left=2mm, right=2mm, top=1mm, bottom=1mm, #1}
\newtcolorbox{takeawaybox}[1][]{colback=VeryLightGray, colframe=RedTitle!55, boxrule=0.35mm,
  arc=1mm, left=2mm, right=2mm, top=1mm, bottom=1mm, #1}
\newtcolorbox{researchbox}[1][]{enhanced, colback=violet!4, colframe=violet!60!black,
  boxrule=0.35mm, arc=1mm, left=2mm, right=2mm, top=1mm, bottom=1mm,
  fonttitle=\bfseries\small, title={Research in the wild}, #1}

\newcommand{\compacttables}{\renewcommand{\arraystretch}{1.15}}
\newcommand{\src}[1]{{\tiny\color{black!55}#1}}

% ---------------------------------------------------------------------------
%% Course map: the recurring "you are here" slide for ECON 282E.
%% Loaded by course_preamble.tex. Pattern follows AI-ECON-2026 lec01_map.tex, but the map
%% shows the whole ten-session course rather than one lecture.
%%
%%   \CourseMap{s}{label}   s = 1..10 highlights session s and prints "you are here: label"
%%                          (e.g. \CourseMap{1}{1.A2}); earlier sessions get a check mark.
%%   \CourseMap{0}{}        full overview, nothing highlighted.
%%
%% Note: TikZ option lists are comma-split before TeX conditionals run, so every \ifnum
%% below sits outside [...] option lists.

\newcommand{\CourseMap}[2]{%
\begin{center}
\begin{tikzpicture}[
    sess/.style={rectangle, rounded corners=2pt, draw=black!45, minimum width=1.34cm,
        minimum height=1.12cm, text width=1.26cm, align=center, inner sep=1pt},
    stage/.style={font=\tiny\bfseries, text=black!70},
]
  \foreach \i/\d/\t/\c in {%
      1/Sep 24/Intro \\ Growth/StageTrunk,
      2/Oct 1/HA $\cdot$ OLG \\ Policy/StageBranch,
      3/Oct 8/Num.\ DP \\ Python/StageMachine,
      4/Oct 15/Numerics \\ I/StageMachine,
      5/Oct 22/Numerics \\ II/StageMachine,
      6/Oct 29/The wall \\ \textbf{Midterm}/StageWall,
      7/Nov 5/ML \\ Deep nets/StageTools,
      8/Nov 12/RL \\ HA $+$ DL/StageTools,
      9/Nov 19/Text \\ LLMs/StageData,
      10/Dec 3/Agents \\ \textbf{Projects}/StageAgents} {
    \node[sess, fill=\c] (n\i) at ({1.46*(\i-1)}, 0)
      {{\scriptsize\bfseries S\i}\\[-1pt]{\tiny\color{black!60}\d}\\[1pt]{\tiny \t}};
    \ifnum\i<#1
      \node[font=\scriptsize, text=green!45!black] at ([xshift=-2.5mm,yshift=-2.2mm]n\i.north east) {$\checkmark$};
    \fi
    \ifnum\i=#1
      \draw[RedTitle, very thick, rounded corners=2pt]
        ([xshift=-0.6mm,yshift=-0.6mm]n\i.south west) rectangle ([xshift=0.6mm,yshift=0.6mm]n\i.north east);
      % keep the label inside the picture at both ends of the row
      \ifnum\i<3
        \node[font=\scriptsize\bfseries, text=RedTitle, anchor=south west] at ([yshift=1.2mm]n\i.north west)
          {$\blacktriangledown$\,you are here: #2};
      \else\ifnum\i>8
        \node[font=\scriptsize\bfseries, text=RedTitle, anchor=south east] at ([yshift=1.2mm]n\i.north east)
          {you are here: #2\,$\blacktriangledown$};
      \else
        \node[font=\scriptsize\bfseries, text=RedTitle, anchor=south] at ([yshift=1.2mm]n\i.north)
          {you are here: #2\,$\blacktriangledown$};
      \fi\fi
    \fi
  }
  % stage band under the sessions
  \foreach \a/\b/\lab/\c in {1/1/trunk/StageTrunk, 2/2/branches/StageBranch,
      3/5/the machine $\cdot$ classical methods/StageMachine, 6/6/the wall/StageWall,
      7/8/new tools/StageTools, 9/9/new data/StageData, 10/10/colleagues/StageAgents} {
    \fill[\c] ([yshift=-1.5mm]n\a.south west) rectangle ([yshift=-4.6mm]n\b.south east);
    \path ([yshift=-1.5mm]n\a.south west) -- ([yshift=-4.6mm]n\b.south east)
      node[midway, stage] {\lab};
  }
  \node[font=\footnotesize\itshape, text=black!70] at ([yshift=-9.5mm]$(n1.south)!0.5!(n10.south)$)
    {Build it on paper $\;\to\;$ solve it on the machine $\;\to\;$ push it to the frontier with AI};
\end{tikzpicture}
\end{center}
}


%%   \decktitle{1}{A1}{Who I Am \& Why This Course}{September 24, 2026}
%%   \begin{document}
%%   \makedecktitle
%%   ...
%%
%% Adapted from Courses/AI-ECON-2026/source/shared_preamble.tex (Metropolis theme,
%% concept/caution/example/takeaway boxes, \sectiondivider). Changes for this course:
%%   * course title block (\decktitle / \makedecktitle) instead of \makebeamertitle
%%   * \zh{...} typesets Chinese (book titles) under pdflatex via CJKutf8
%%   * researchbox for the "Research in the wild" frames
%%   * section pages off; TikZ libraries and the course map loaded here
%% Build with pdflatex only (two passes); tools/build_slides.py does this for you.

\documentclass[10pt,english,aspectratio=169]{beamer}

\usepackage[T1]{fontenc}
\usepackage{textcomp}
\usepackage[utf8]{inputenc}
\usepackage{babel}
\usepackage{CJKutf8}

\usepackage{amsmath, amssymb, amsthm, graphicx}
\usepackage{booktabs, multirow, tabularx}
\usepackage{tikz}
\usetikzlibrary{positioning,arrows.meta,shapes,calc,fit,backgrounds}
\usepackage{pgfplots}
\pgfplotsset{compat=1.16}
\usepackage[most]{tcolorbox}
\tcbuselibrary{skins, breakable}
\usepackage{listings}
\usepackage{xcolor}

\lstset{
  basicstyle=\ttfamily\small,
  keywordstyle=\color{blue!80!black}\bfseries,
  commentstyle=\color{green!50!black}\itshape,
  stringstyle=\color{orange!80!black},
  backgroundcolor=\color{gray!8},
  frame=single,
  rulecolor=\color{gray!40},
  breaklines=true,
  showstringspaces=false,
  numbers=none,
}

\usetheme{metropolis}
\usefonttheme{professionalfonts}
\metroset{sectionpage=none}

\definecolor{LightGray}{RGB}{230,230,230}
\definecolor{RedTitle}{RGB}{0,0,200}
\definecolor{VeryLightGray}{RGB}{245,245,245}
\definecolor{DarkText}{RGB}{0,0,0}
\definecolor{UCRBlue}{RGB}{0,45,106}
\definecolor{UCRGold}{RGB}{241,171,0}

% Stage colours of the course arc (used by the course map and elsewhere)
\colorlet{StageTrunk}{blue!12}
\colorlet{StageBranch}{cyan!14}
\colorlet{StageMachine}{gray!18}
\colorlet{StageWall}{red!14}
\colorlet{StageTools}{orange!20}
\colorlet{StageData}{green!16}
\colorlet{StageAgents}{violet!14}

\hypersetup{bookmarks=false,breaklinks=true,pdfborder={0 0 0},colorlinks=true,
  linkcolor=DarkText,urlcolor=RedTitle}

\setbeamercolor{background canvas}{bg=white}
\setbeamercolor{title}{fg=RedTitle}
\setbeamercolor{frametitle}{bg=VeryLightGray,fg=DarkText}
\setbeamercolor{normal text}{fg=DarkText}
\setbeamercolor{alerted text}{fg=RedTitle}
\setbeamersize{text margin left=5mm, text margin right=5mm}
\addtobeamertemplate{frametitle}{}{\vspace{-0.5em}}

\setbeamertemplate{navigation symbols}{}
\setbeamertemplate{footline}{%
  \hspace{5mm}{\usebeamerfont{page number in head/foot}\color{black!45}%
  ECON 282E \,$\cdot$\, \insertshortdeck}\hfill%
  \usebeamercolor[fg]{page number in head/foot}%
  \usebeamerfont{page number in head/foot}%
  \insertframenumber\,/\,\inserttotalframenumber\kern1em\vskip2pt%
}

% ---------------------------------------------------------------------------
% Chinese text under pdflatex (book titles etc.)
\newcommand{\zh}[1]{\begin{CJK*}{UTF8}{gbsn}#1\end{CJK*}}

% ---------------------------------------------------------------------------
% Deck title block
%   \decktitle{session}{deck id}{title}{date}
\newcommand{\insertshortdeck}{}
\newcommand{\decksession}{}
\newcommand{\deckid}{}
\newcommand{\decktitletext}{}
\newcommand{\deckdate}{}
\newcommand{\decktitle}[4]{%
  \renewcommand{\decksession}{#1}%
  \renewcommand{\deckid}{#2}%
  \renewcommand{\decktitletext}{#3}%
  \renewcommand{\deckdate}{#4}%
  \renewcommand{\insertshortdeck}{Session #1 \,$\cdot$\, Deck #1.#2}%
  \title{#3}%
  \author{Zhigang Feng}%
  \date{#4}%
}
\newcommand{\makedecktitle}{%
  \begin{frame}[plain,noframenumbering]
    \begin{tikzpicture}[remember picture,overlay]
      \fill[UCRBlue] (current page.north west) rectangle ([yshift=-0.55cm]current page.north east);
      \fill[UCRGold] ([yshift=-0.55cm]current page.north west) rectangle ([yshift=-0.62cm]current page.north east);
      \node[anchor=west, text=white, font=\small\bfseries] at ([xshift=5mm,yshift=-0.275cm]current page.north west)
        {ECON 282E \,$\cdot$\, Foundations of Macroeconomics};
      \node[anchor=east, text=white, font=\small] at ([xshift=-5mm,yshift=-0.275cm]current page.north east)
        {UC Riverside \,$\cdot$\, Fall 2026};
    \end{tikzpicture}
    \vspace*{1.1cm}
    {\usebeamerfont{title}\color{black!55}\large Session \decksession\ \,$\cdot$\, Deck \decksession.\deckid\par}
    \vspace{0.25cm}
    {\usebeamerfont{title}\color{RedTitle}\LARGE\bfseries \decktitletext\par}
    \vspace{0.7cm}
    {\large Zhigang Feng\par}
    \vspace{0.15cm}
    {\color{black!65}\deckdate\par}
    \vfill
    {\footnotesize\itshape\color{black!55}Build it on paper $\;\to\;$ solve it on the machine $\;\to\;$ push it to the frontier with AI\par}
  \end{frame}}

% ---------------------------------------------------------------------------
% Section divider (plain frame, not counted)
\newcommand{\sectiondivider}[2]{%
\begin{frame}[plain,noframenumbering]
\begin{center}
\vfill
{\Large\textbf{#1}\par}
\vspace{0.35cm}
{\normalsize #2\par}
\vfill
\end{center}
\end{frame}
}

% ---------------------------------------------------------------------------
% Boxes
\newtcolorbox{conceptbox}[1][]{colback=blue!3, colframe=RedTitle!55, boxrule=0.35mm,
  arc=1mm, left=2mm, right=2mm, top=1mm, bottom=1mm, #1}
\newtcolorbox{cautionbox}[1][]{colback=red!3, colframe=red!50!black, boxrule=0.35mm,
  arc=1mm, left=2mm, right=2mm, top=1mm, bottom=1mm, #1}
\newtcolorbox{examplebox}[1][]{colback=green!3, colframe=green!45!black, boxrule=0.35mm,
  arc=1mm, left=2mm, right=2mm, top=1mm, bottom=1mm, #1}
\newtcolorbox{takeawaybox}[1][]{colback=VeryLightGray, colframe=RedTitle!55, boxrule=0.35mm,
  arc=1mm, left=2mm, right=2mm, top=1mm, bottom=1mm, #1}
\newtcolorbox{researchbox}[1][]{enhanced, colback=violet!4, colframe=violet!60!black,
  boxrule=0.35mm, arc=1mm, left=2mm, right=2mm, top=1mm, bottom=1mm,
  fonttitle=\bfseries\small, title={Research in the wild}, #1}

\newcommand{\compacttables}{\renewcommand{\arraystretch}{1.15}}
\newcommand{\src}[1]{{\tiny\color{black!55}#1}}

% ---------------------------------------------------------------------------
%% Course map: the recurring "you are here" slide for ECON 282E.
%% Loaded by course_preamble.tex. Pattern follows AI-ECON-2026 lec01_map.tex, but the map
%% shows the whole ten-session course rather than one lecture.
%%
%%   \CourseMap{s}{label}   s = 1..10 highlights session s and prints "you are here: label"
%%                          (e.g. \CourseMap{1}{1.A2}); earlier sessions get a check mark.
%%   \CourseMap{0}{}        full overview, nothing highlighted.
%%
%% Note: TikZ option lists are comma-split before TeX conditionals run, so every \ifnum
%% below sits outside [...] option lists.

\newcommand{\CourseMap}[2]{%
\begin{center}
\begin{tikzpicture}[
    sess/.style={rectangle, rounded corners=2pt, draw=black!45, minimum width=1.34cm,
        minimum height=1.12cm, text width=1.26cm, align=center, inner sep=1pt},
    stage/.style={font=\tiny\bfseries, text=black!70},
]
  \foreach \i/\d/\t/\c in {%
      1/Sep 24/Intro \\ Growth/StageTrunk,
      2/Oct 1/HA $\cdot$ OLG \\ Policy/StageBranch,
      3/Oct 8/Num.\ DP \\ Python/StageMachine,
      4/Oct 15/Numerics \\ I/StageMachine,
      5/Oct 22/Numerics \\ II/StageMachine,
      6/Oct 29/The wall \\ \textbf{Midterm}/StageWall,
      7/Nov 5/ML \\ Deep nets/StageTools,
      8/Nov 12/RL \\ HA $+$ DL/StageTools,
      9/Nov 19/Text \\ LLMs/StageData,
      10/Dec 3/Agents \\ \textbf{Projects}/StageAgents} {
    \node[sess, fill=\c] (n\i) at ({1.46*(\i-1)}, 0)
      {{\scriptsize\bfseries S\i}\\[-1pt]{\tiny\color{black!60}\d}\\[1pt]{\tiny \t}};
    \ifnum\i<#1
      \node[font=\scriptsize, text=green!45!black] at ([xshift=-2.5mm,yshift=-2.2mm]n\i.north east) {$\checkmark$};
    \fi
    \ifnum\i=#1
      \draw[RedTitle, very thick, rounded corners=2pt]
        ([xshift=-0.6mm,yshift=-0.6mm]n\i.south west) rectangle ([xshift=0.6mm,yshift=0.6mm]n\i.north east);
      % keep the label inside the picture at both ends of the row
      \ifnum\i<3
        \node[font=\scriptsize\bfseries, text=RedTitle, anchor=south west] at ([yshift=1.2mm]n\i.north west)
          {$\blacktriangledown$\,you are here: #2};
      \else\ifnum\i>8
        \node[font=\scriptsize\bfseries, text=RedTitle, anchor=south east] at ([yshift=1.2mm]n\i.north east)
          {you are here: #2\,$\blacktriangledown$};
      \else
        \node[font=\scriptsize\bfseries, text=RedTitle, anchor=south] at ([yshift=1.2mm]n\i.north)
          {you are here: #2\,$\blacktriangledown$};
      \fi\fi
    \fi
  }
  % stage band under the sessions
  \foreach \a/\b/\lab/\c in {1/1/trunk/StageTrunk, 2/2/branches/StageBranch,
      3/5/the machine $\cdot$ classical methods/StageMachine, 6/6/the wall/StageWall,
      7/8/new tools/StageTools, 9/9/new data/StageData, 10/10/colleagues/StageAgents} {
    \fill[\c] ([yshift=-1.5mm]n\a.south west) rectangle ([yshift=-4.6mm]n\b.south east);
    \path ([yshift=-1.5mm]n\a.south west) -- ([yshift=-4.6mm]n\b.south east)
      node[midway, stage] {\lab};
  }
  \node[font=\footnotesize\itshape, text=black!70] at ([yshift=-9.5mm]$(n1.south)!0.5!(n10.south)$)
    {Build it on paper $\;\to\;$ solve it on the machine $\;\to\;$ push it to the frontier with AI};
\end{tikzpicture}
\end{center}
}


\graphicspath{{./figures/}}

\lstset{language=Python, basicstyle=\ttfamily\scriptsize, frame=none,
        backgroundcolor=\color{gray!8}, aboveskip=2pt, belowskip=2pt}

\decktitle{5}{A2}{The Model, and the Parameter That Scales Risk}{Thursday, October 22, 2026}

\begin{document}

\makedecktitle

%=============================================================================
\begin{frame}{Where We Are}
\CourseMap{5}{5.A2}
\vspace{-1mm}
\begin{center}
\small \textbf{Block A:} 5.A1 what perturbation computes $\;\cdot\;$ \textbf{5.A2} the model and the parameter $\;\cdot\;$ 5.A3 first order $\;\cdot\;$ 5.A4 linear RE systems $\;\cdot\;$ 5.A5 higher order, risk, pruning $\;\cdot\;$ 5.A6 the worked example.
\end{center}
\end{frame}

%=============================================================================
\begin{frame}{The Model We Will Expand}
\begin{columns}[T]
\begin{column}{0.55\textwidth}
\small
\vspace{-2mm}
\begin{align*}
\max\;&\mathbb{E}_{0}\sum_{t=0}^{\infty}\beta^{t}\ln c_{t}\\
\text{s.t.}\;\; & c_{t}+k_{t+1}=e^{z_{t}}k_{t}^{\alpha}\\
 & z_{t}=\rho z_{t-1}+\sigma\varepsilon_{t},\quad \varepsilon_t\sim N(0,1)
\end{align*}
Equilibrium conditions:
\begin{align*}
\frac{1}{c_{t}} &=\beta\,\mathbb{E}_{t}\!\left[\frac{1}{c_{t+1}}\,\alpha e^{z_{t+1}}k_{t+1}^{\alpha-1}\right]\\
c_{t}+k_{t+1} &= e^{z_{t}}k_{t}^{\alpha}
\end{align*}
\textbf{States} $(k_t,z_t)$. \textbf{Controls} $(c_t,k_{t+1})$. The object we want is the pair of \textbf{policy functions} $c(k,z)$ and $k'(k,z)$.
\end{column}
\begin{column}{0.42\textwidth}
\begin{conceptbox}[title=\textbf{Why this model}]\small
Log utility with \textbf{full depreciation} is the Brock--Mirman benchmark of 1.B1 --- the one model in this course whose policy is known in closed form. Every method gets tested here first.
\end{conceptbox}
\medskip
\begin{examplebox}\footnotesize
$\delta=1$ is not realistic. It is diagnostic. 5.A6 moves to the quarterly calibration with $\delta=0.025$, where no closed form exists.
\end{examplebox}
\end{column}
\end{columns}
\vspace{1mm}
\src{Quant\_Macro Lec.~9, ``Stochastic neoclassical growth model''. Calibration $\alpha=0.36$, $\beta=0.95$, as in S3.}
\end{frame}

%=============================================================================
\begin{frame}{The One Thing We Know Exactly}
\begin{columns}[T]
\begin{column}{0.52\textwidth}
\small
Guess and verify gives the exact policy:
\begin{align*}
c_{t} &= (1-\alpha\beta)\,e^{z_{t}}k_{t}^{\alpha}\\
k_{t+1} &= \alpha\beta\,e^{z_{t}}k_{t}^{\alpha}
\end{align*}
and therefore the deterministic steady state
\begin{align*}
k^{*} &= (\alpha\beta)^{\frac{1}{1-\alpha}} = 0.1870\\
c^{*} &= (\alpha\beta)^{\frac{\alpha}{1-\alpha}}-(\alpha\beta)^{\frac{1}{1-\alpha}}\\
z^{*} &= 0 .
\end{align*}
The steady state is the point we expand around; the closed-form policy is the answer we will grade against.
\end{column}
\begin{column}{0.45\textwidth}
\begin{takeawaybox}\footnotesize
\textbf{Two different objects.} Almost every model has a computable steady state. Almost none has a closed-form \emph{policy}. Perturbation needs only the first --- which is why it generalizes.
\end{takeawaybox}
\medskip
\begin{examplebox}\footnotesize
1.B1: ``If a method cannot recover $k'=\alpha\beta z k^{\alpha}$, do not trust it on anything harder.'' VFI was tested here in S3. Perturbation is tested here in three slides.
\end{examplebox}
\end{column}
\end{columns}
\vspace{1mm}
\src{Quant\_Macro Lec.~9, ``Solution and steady state''. Ledger \texttt{brock\_mirman}: $k^{*}=0.187032$.}
\end{frame}

%=============================================================================
\begin{frame}{The Goal, Written as Functional Equations}
\small
Substituting the unknown policies into the equilibrium conditions gives a system in which the unknowns are \emph{functions}:
\begin{align*}
\frac{1}{c(k_{t},z_{t})} &=\beta\,\mathbb{E}_{t}\!\left[\frac{\alpha\,e^{\rho z_{t}+\sigma\varepsilon_{t+1}}\,k(k_{t},z_{t})^{\alpha-1}}
{c\big(k(k_{t},z_{t}),\,\rho z_{t}+\sigma\varepsilon_{t+1}\big)}\right]\\[1mm]
c(k_{t},z_{t})+k(k_{t},z_{t}) &= e^{z_{t}}k_{t}^{\alpha}
\end{align*}
\medskip
\begin{columns}[T]
\begin{column}{0.5\textwidth}
\begin{conceptbox}[title=\textbf{Read the composition}]\footnotesize
$c$ appears \emph{evaluated at next period's state}, which is itself the unknown $k(\cdot)$. That nesting is why this is a functional equation and not an algebraic one.
\end{conceptbox}
\end{column}
\begin{column}{0.47\textwidth}
\begin{examplebox}\footnotesize
Once $c(\cdot)$ and $k(\cdot)$ are known, everything follows: simulate, compute moments, produce impulse responses, do welfare.
\end{examplebox}
\end{column}
\end{columns}
\vspace{1mm}
\src{Quant\_Macro Lec.~9, ``The goal''.}
\end{frame}

%=============================================================================
\begin{frame}{A Parameter That Scales the Risk}
\begin{columns}[T]
\begin{column}{0.55\textwidth}
\small
There is a problem with expanding in $(k,z)$ alone: we want to know what \textbf{uncertainty itself} does, and uncertainty is not a state.
\medskip

The device is to introduce an artificial parameter $\chi$ multiplying the shock:
\[
  z_{t}=\rho z_{t-1}+\chi\,\sigma\varepsilon_{t}
\]
\begin{itemize}\small
  \item $\chi=1$: the model we care about.
  \item $\chi=0$: the \textbf{deterministic} economy, where we know how to compute everything.
\end{itemize}
\medskip
The policies become $c(k_t,z_t;\chi)$ and $k(k_t,z_t;\chi)$, and $\chi$ is treated as a third argument to differentiate with respect to.
\end{column}
\begin{column}{0.42\textwidth}
\begin{cautionbox}[title=\textbf{Notation}]\footnotesize
The literature, including Fern\'andez-Villaverde \& Guerr\'on-Quintana, writes $\lambda$ for this parameter. This course reserves $\lambda$ for the \textbf{invariant distribution} (2.A3), so we write $\chi$. Nothing else changes.
\end{cautionbox}
\medskip
\begin{conceptbox}[title=\textbf{The move}]\footnotesize
Risk has been turned into a \emph{parameter we can differentiate}. That is the entire trick, and everything in 5.A5 is a consequence of it.
\end{conceptbox}
\end{column}
\end{columns}
\vspace{1mm}
\src{Quant\_Macro Lec.~9, ``A perturbation solution''.}
\end{frame}

%=============================================================================
\begin{frame}{Why Expand Around $\chi=0$?}
\begin{columns}[T]
\begin{column}{0.55\textwidth}
\small
Because $\chi=0$ is the one point where the answer is available: no uncertainty, a deterministic steady state, and policies we can compute directly.
\medskip

Expanding around it measures \textbf{how a small amount of risk changes behaviour relative to the deterministic economy} --- which is exactly the economic question.
\medskip

The by-products are immediate:
\begin{itemize}\small
  \item first order gives a \textbf{linear state-space system}, so impulse responses are one matrix multiplication;
  \item second-order moments follow in closed form, with no simulation.
\end{itemize}
\end{column}
\begin{column}{0.42\textwidth}
\begin{cautionbox}[title=\textbf{And the price}]\small
Accuracy deteriorates with large shocks, strong nonlinearity, and distance from the steady state. We do not have to take that on trust --- 5.A6 measures exactly where it happens.
\end{cautionbox}
\medskip
\begin{examplebox}\footnotesize
Note what is being assumed: that the solution is \emph{smooth in $\chi$ at $\chi=0$}. Models with occasionally binding constraints are not.
\end{examplebox}
\end{column}
\end{columns}
\vspace{1mm}
\src{Quant\_Macro Lec.~9, ``Why Perturb Around $\lambda=0$?''.}
\end{frame}

%=============================================================================
\begin{frame}{What Exactly Gets Differentiated}
\begin{columns}[T]
\begin{column}{0.54\textwidth}
\small
We build a local approximation around $(k^{*},0;0)$. Taking the equilibrium conditions
\begin{align*}
\frac{1}{c(k_{t},z_{t};\chi)} &=\beta\mathbb{E}_{t}\!\left[\cdots\right]\\
c(k_{t},z_{t};\chi)+k(k_{t},z_{t};\chi) &= e^{z_{t}}k_{t}^{\alpha}
\end{align*}
we differentiate with respect to $k_t$, $z_t$ and $\chi$, and evaluate at $(k^{*},0;0)$.
\medskip

\textbf{The unknowns become coefficients.} $c_k$, $c_z$, $c_\chi$, $k_k$, $k_z$, $k_\chi$ are numbers, not functions --- the partial derivatives of the true policies at one point.
\end{column}
\begin{column}{0.43\textwidth}
\begin{conceptbox}[title=\textbf{Taylor's theorem, twice over}]\small
Once to justify writing the policy as a series at all, and once to say what the truncation error is. Both need the smoothness the IFT gave us in 5.A1.
\end{conceptbox}
\medskip
\begin{takeawaybox}\footnotesize
``Solve the model'' has become ``compute a few derivatives at a point we already know''.
\end{takeawaybox}
\end{column}
\end{columns}
\vspace{1mm}
\src{Quant\_Macro Lec.~9, ``Taylor's theorem'' --- including the two frames of that title carried in the predecessor deck and dropped from the current one.}
\end{frame}

%=============================================================================
\begin{frame}{The Expansion, Written Out}
\small
To second order, around $(k^{*},0;0)$, writing $\hat{k}_t=k_t-k^{*}$:
\begin{align*}
c_t \;=\;& c^{*} + c_k\hat{k}_t + c_z z_t + c_\chi \chi \\
 &+ \tfrac12\Big(c_{kk}\hat{k}_t^{2} + 2c_{kz}\hat{k}_t z_t + c_{zz}z_t^{2}
   + 2c_{k\chi}\hat{k}_t\chi + 2c_{z\chi}z_t\chi + c_{\chi\chi}\chi^{2}\Big)+\cdots\\[1mm]
k_{t+1} \;=\;& k^{*} + k_k\hat{k}_t + k_z z_t + k_\chi \chi \\
 &+ \tfrac12\Big(k_{kk}\hat{k}_t^{2} + 2k_{kz}\hat{k}_t z_t + k_{zz}z_t^{2}
   + 2k_{k\chi}\hat{k}_t\chi + 2k_{z\chi}z_t\chi + k_{\chi\chi}\chi^{2}\Big)+\cdots
\end{align*}
\medskip
\begin{columns}[T]
\begin{column}{0.48\textwidth}
\begin{conceptbox}[title=\textbf{First order}]\footnotesize
Keep the linear terms only: six unknowns, $\{c_k,c_z,c_\chi,k_k,k_z,k_\chi\}$. 5.A3 shows that two of them are zero.
\end{conceptbox}
\end{column}
\begin{column}{0.48\textwidth}
\begin{examplebox}\footnotesize
$c_{\chi\chi}$ is the term that will carry \textbf{precautionary behaviour}. It is a second-order object, which is already the answer to ``why go beyond first order''.
\end{examplebox}
\end{column}
\end{columns}
\vspace{1mm}
\src{Quant\_Macro Lec.~9, ``Taylor Expansion of Policy Functions''.}
\end{frame}

%=============================================================================
\begin{frame}{Measured: On This Model, the Linear Solution Is Not an Approximation}
\begin{columns}[T]
\begin{column}{0.52\textwidth}
\small
Take logs of the exact policy:
\[
  \ln k_{t+1}=\ln(\alpha\beta)+z_t+\alpha\ln k_t ,
\]
and since $\ln k^{*}=\ln(\alpha\beta)/(1-\alpha)$,
\[
  \widehat{\ln k}_{t+1}=\alpha\,\widehat{\ln k}_{t}+z_t .
\]
\textbf{The true policy is exactly linear in logs.} So a first-order perturbation in logs should return $\alpha$ and $1$ --- not approximately.
\medskip

Solved numerically, it does:
\end{column}
\begin{column}{0.45\textwidth}
\begin{center}\footnotesize
\begin{tabularx}{\linewidth}{@{}>{\raggedright\arraybackslash}X r@{}}
\toprule
\textbf{Measured} & \\
\midrule
$k$ on $\hat{k}$, $\;k$ on $z$ & $0.36$,\; $1.00$ \\
$c$ on $\hat{k}$, $\;c$ on $z$ & $0.36$,\; $1.00$ \\
deviation from exact & $1.1\times10^{-16}$ \\
policy error, \emph{in levels} & $4.4\times10^{-16}$ \\
every second-order term & $<10^{-9}$ \\
\bottomrule
\end{tabularx}
\end{center}
\medskip
\begin{takeawaybox}\footnotesize
Machine precision, over $k\in[0.25k^{*},2.5k^{*}]$ --- a tenfold range, not a neighbourhood.
\end{takeawaybox}
\end{column}
\end{columns}
\vspace{1mm}
\src{Ledger \texttt{brock\_mirman}. Coefficients by generalized Schur; the second-order terms come from an independent residual-matching solve and vanish to the precision of its finite differences.}
\end{frame}

%=============================================================================
\begin{frame}{What That Benchmark Can and Cannot Tell You}
\begin{columns}[T]
\begin{column}{0.53\textwidth}
\small
\textbf{What it certifies.} The algebra, the code, the steady state, the sign conventions, the transposes. If perturbation cannot return $0.36$ and $1.00$ here, the implementation is wrong and nothing else matters.
\medskip

\textbf{What it cannot certify.} Anything about accuracy. The error is zero \emph{because the model is log-linear}, not because perturbation is good. A benchmark that a method solves exactly measures nothing about the method.
\medskip

This is the same trap as 3.A2's beautifully converged value function with a bad policy: the diagnostic that looks most reassuring is the one carrying least information.
\end{column}
\begin{column}{0.44\textwidth}
\begin{cautionbox}[title=\textbf{Why $\delta=1$ does this}]\small
Full depreciation plus log utility makes the saving \emph{rate} constant at $\alpha\beta$. Nothing is left for curvature to act on, so certainty equivalence is exact and every higher-order term is zero.
\end{cautionbox}
\medskip
\begin{examplebox}\footnotesize
Hence 5.A6 runs the quarterly model with $\delta=0.025$, where the second-order terms are not zero and the error is not zero either.
\end{examplebox}
\end{column}
\end{columns}
\vspace{1mm}
\src{Ledger \texttt{brock\_mirman}, field \texttt{note}.}
\end{frame}

%=============================================================================
\begin{frame}{Takeaways}
\begin{takeawaybox}
\begin{itemize}\small
  \item The model is the \textbf{Brock--Mirman benchmark} of 1.B1: log utility, $\delta=1$, closed-form policy, steady state $k^{*}=0.1870$.
  \item Perturbation needs only the \textbf{steady state}, not the closed form --- which is exactly why it survives into models that have no closed form.
  \item Risk becomes differentiable by introducing $\chi$ with $z'=\rho z+\chi\sigma\varepsilon'$: $\chi=0$ is the deterministic economy, $\chi=1$ the one we want.
  \item ``Solve the model'' has become ``find $\{c_k,c_z,c_\chi,k_k,k_z,k_\chi\}$, then the second-order terms, then\dots''.
  \item Measured: on this model the first-order solution reproduces the exact policy to \textbf{$10^{-16}$ over a tenfold range of capital}, and every second-order term is zero.
  \item Which is a \textbf{correctness} test, not an accuracy test. A benchmark solved exactly tells you nothing about how the method behaves where it is actually needed.
\end{itemize}
\end{takeawaybox}
\end{frame}

%=============================================================================
\begin{frame}{Bridge: Where the Equations Come From}
\begin{columns}[T]
\begin{column}{0.53\textwidth}
\small
We now have unknowns --- six of them at first order --- and no equations for them yet. The next deck writes the equilibrium conditions as a single object $F(k,z;\chi)=0$, observes that it is an \textbf{identity}, and differentiates it.
\medskip

Two results fall out that are worth waiting for: two of the six unknowns turn out to be \textbf{exactly zero}, and the remaining system is exactly square.
\end{column}
\begin{column}{0.44\textwidth}
\begin{conceptbox}[title=\textbf{Next (5.A3)}]\small
$F\equiv 0 \Rightarrow F_k=F_z=F_\chi=0$; six unknowns, four equations, and why that is not a mistake. Certainty equivalence, and what it costs.
\end{conceptbox}
\medskip
\begin{examplebox}\footnotesize
Keep the closed form in view. Everything 5.A3 derives can be checked against $\alpha$ and $1$.
\end{examplebox}
\end{column}
\end{columns}
\end{frame}

\end{document}
